% BHU PYQ -- Commerce: Business Economics (2025-26 NEP)
\documentclass[11pt,a4paper]{article}
\usepackage[a4paper,top=2.5cm,bottom=2.5cm,left=2.8cm,right=2.8cm,headheight=14pt]{geometry}
\usepackage{amsmath,amssymb}
\usepackage{fontspec}
\setmainfont{Kohinoor Devanagari}
\usepackage{microtype,fancyhdr,enumitem,hyperref,lastpage,parskip}

\hypersetup{colorlinks=true,urlcolor=black,linkcolor=black,pdftitle={COMMN11: Business Economics -- BHU NEP 2025-26}}

\pagestyle{fancy}\fancyhf{}
\renewcommand{\headrulewidth}{0.4pt}\renewcommand{\footrulewidth}{0.4pt}
\fancyhead[L]{\small \textit{Banaras Hindu University}}
\fancyhead[R]{\small \textit{COMMN11: Business Economics (2025-26)}}
\fancyfoot[C]{\small Page \thepage\ of \pageref{LastPage}}

\newlist{parts}{enumerate}{2}
\setlist[parts,1]{label=(\alph*),leftmargin=*,itemsep=4pt,topsep=3pt}
\setlist[parts,2]{label=(\roman*),leftmargin=*,itemsep=2pt,topsep=2pt}

\newcommand{\pts}[1]{\hfill \textit{[#1]}}

\begin{document}

\begin{center}
  {\large\bfseries BANARAS HINDU UNIVERSITY}\\[2pt]
  {\normalsize B. Com. (Hons.) Semester I Examination, 2025-26 (NEP)}\\[6pt]
  \rule{0.8\linewidth}{0.5pt}\\[6pt]
  {\large\bfseries COMMERCE}\\[4pt]
  {\large\bfseries Paper Code: COMMN11 : Business Economics}\\[4pt]
  {\normalsize Time: 2:30 Hours \hfill Full Marks: 70}\\[2pt]
  {\small REF NO. CECONF/2025-26/PS/25}
\end{center}
\rule{\linewidth}{0.4pt}
\smallskip

\noindent \textit{Instructions: Answer \textbf{all} questions. All questions carry equal marks (14 marks each). Write your Roll No.\ at the top immediately on receipt of this question paper.}

\smallskip
\noindent \textit{सभी प्रश्नों के उत्तर दीजिए। सभी प्रश्नों के अंक समान हैं।}

\bigskip

\begin{enumerate}[label=\textbf{\arabic*.},leftmargin=*,itemsep=14pt]

  \item Explain the concept and law of demand. Why does a demand curve generally slope downward from left to right? Discuss various exceptions of this law.\pts{14}
  
  माँग की अवधारणा तथा नियम की व्याख्या कीजिए। माँग वक्र सामान्यतः बाएँ से दाएँ नीचे की ओर क्यों झुकता है? इस नियम के विभिन्न अपवादों की चर्चा कीजिए।

  \begin{center}\textbf{OR / अथवा}\end{center}

  Define the price and cross elasticity of demand and explain its various degrees. Discuss the main methods of measuring the elasticity of demand.\pts{14}
  
  माँग की कीमत तथा आड़ी लोच की परिभाषा दीजिए तथा इसकी विभिन्न श्रेणियों की व्याख्या कीजिए। माँग की लोच को मापने की मुख्य विधियों की चर्चा कीजिए।

  \item ``There are many advantages available to large-scale production, yet it is not free from limitations.'' In the light of this statement, briefly discuss various economies and diseconomies of scale.\pts{14}
  
  ``बड़े पैमाने के उत्पादन में अनेक लाभ उपलब्ध हैं, फिर भी यह सीमाओं से मुक्त नहीं है।'' इस कथन के प्रकाश में, पैमाने की विभिन्न आंतरिक तथा बाह्य मितव्ययिताओं एवं अमितव्ययिताओं की संक्षेप में चर्चा कीजिए।

  \begin{center}\textbf{OR / अथवा}\end{center}

  Briefly describe the cost-output relationships in different time periods.\pts{14}
  
  विभिन्न समयावधियों में लागत-उत्पादन सम्बन्धों का संक्षेप में वर्णन कीजिए।

  \item How does a firm attain equilibrium under condition of perfect competition in short- and long-run? Explain with the help of suitable diagrams.\pts{14}
  
  एक फर्म पूर्ण प्रतियोगिता की स्थिति में अल्पकाल तथा दीर्घकाल में साम्य (संतुलन) कैसे प्राप्त करती है? उपयुक्त चित्रों की सहायता से समझाइए।

  \begin{center}\textbf{OR / अथवा}\end{center}

  Discuss the main features of Monopoly. How are price and output determined under Monopoly?\pts{14}
  
  एकाधिकार की मुख्य विशेषताओं की चर्चा कीजिए। एकाधिकार के अंतर्गत कीमत तथा उत्पादन का निर्धारण किस प्रकार होता है?

  \item What is Price Discrimination? Under what conditions is Price Discrimination possible and profitable?\pts{14}
  
  कीमत विभेद क्या है? किन परिस्थितियों में कीमत विभेद संभव तथा लाभदायक होता है?

  \begin{center}\textbf{OR / अथवा}\end{center}

  Discuss the Marginal Productivity Theory of Distribution. What are its main criticisms?\pts{14}
  
  वितरण के सीमांत उत्पादकता सिद्धान्त की चर्चा कीजिए। इसकी मुख्य आलोचनाएँ क्या हैं?

  \item Write short notes on any \textbf{two} of the following:\pts{2 $\times$ 7 = 14}
  
  निम्नलिखित में से किन्हीं दो पर संक्षिप्त टिप्पणियाँ लिखिए:
  \begin{parts}
    \item Indifference Curve Analysis (उदासीनता वक्र विश्लेषण)
    \item Returns to Scale (पैमाने के प्रतिफल)
    \item Kinked Demand Curve (विकुंचित माँग वक्र)
    \item Ricardian Theory of Rent (रिकार्डो का लगान सिद्धान्त)
  \end{parts}

\end{enumerate}

\end{document}
